% !TeX program = xelatex
% Save this file with UTF-8 encoding.

\documentclass[a4paper,10pt]{article}

\usepackage[
  a4paper,
  top=26mm,
  bottom=12mm,
  left=25mm,
  right=25mm
]{geometry}

% Fonts and Chinese support
\usepackage{fontspec}
\usepackage{xeCJK}

\setmainfont{TeX Gyre Termes}
\setsansfont{TeX Gyre Heros}

\setCJKmainfont{FandolSong-Regular}[
  BoldFont=FandolSong-Bold,
  ItalicFont=FandolKai-Regular
]
\setCJKsansfont{FandolHei-Regular}[
  BoldFont=FandolHei-Bold
]

\usepackage{graphicx}
\usepackage{amsmath,amssymb}
\usepackage{anyfontsize}
\usepackage{xcolor}
\usepackage{array,tabularx}
\usepackage{microtype}

\definecolor{DLADEBlue}{HTML}{003B70}
\definecolor{DLADEAccent}{HTML}{0066C5}
\definecolor{DLADEGold}{HTML}{FFCA05}
\definecolor{DLADEMuted}{HTML}{59758E}

\definecolor{Highlight}{HTML}{D7191C}

\usepackage[
  colorlinks=true,
  urlcolor=DLADEAccent,
  linkcolor=DLADEAccent
]{hyperref}

\pagestyle{empty}
\setlength{\parindent}{0pt}
\setlength{\parskip}{0pt}
\renewcommand{\arraystretch}{1.0}
\emergencystretch=1em

% -----------------------------------------------------------------------------
% EDITABLE SEMINAR DETAILS
% -----------------------------------------------------------------------------

\newcommand{\SeminarSeries}{DDES - DUT Differential Equations Seminar}
\newcommand{\SeminarTitle}{The linear stability of relative equilibria in the N-body problem }
\newcommand{\SpeakerName}{JieXin Sun}
\newcommand{\SpeakerAffiliation}{孙洁昕, PostDoc@DUT}
%\textcolor{Highlight}{\textbf{111-A}}
\newcommand{\SeminarVenue}{Seminar Room 114}
%\textcolor{Highlight}{2:00pm - 2:45pm}
\newcommand{\SeminarTime}{9:00am - 9:45am, Thursday, October 8, 2026}
\newcommand{\GroupURL}{https://dutdynamics.github.io/ddes/}

\hypersetup{
  pdftitle={\SeminarTitle},
  pdfauthor={DLADE Group, School of Mathematical Sciences, Dalian University of Technology},
  pdfsubject={DLADE seminar report}
}

\begin{document}

% =============================================================================
% UPPER HALF: IDENTITY, LOGISTICS, TITLE AND ABSTRACT
% =============================================================================

\noindent
\begin{minipage}[t][\dimexpr0.50\textheight+35mm\relax][t]{\textwidth}
\vspace{0pt}

% -----------------------------------------------------------------------------
% UNIVERSITY IDENTITY
% -----------------------------------------------------------------------------

\begin{tabularx}{\textwidth}{@{}X@{\hspace{8mm}}p{62mm}@{}}
  \begin{minipage}[t]{\linewidth}
    \vspace{0pt}
    \includegraphics[
      trim=66 45 60 55,
      clip,
      width=79mm
    ]{DUT_symbol.pdf}
  \end{minipage}
  &
  \begin{minipage}[t]{\linewidth}
    \vspace{0pt}
    \hfill
    \begin{minipage}[t]{60mm}
      \raggedleft
      {\sffamily\bfseries\fontsize{9.5}{11.2}\selectfont
        School of Mathematical Sciences\par
        Dalian University of Technology\par
      }

      \vspace{1mm}

      {\sffamily\fontsize{8.4}{10}\selectfont
        Linggong Road 2\par
        116024 Dalian, P.R. China\par
      }

      \vspace{1mm}

      {\sffamily\fontsize{6.2}{7.4}\selectfont\color{DLADEAccent}
        \href{\GroupURL}{dutdynamics.github.io/ddes}\par
      }
    \end{minipage}
  \end{minipage}
\end{tabularx}

% -----------------------------------------------------------------------------
% SEMINAR LOGISTICS
% -----------------------------------------------------------------------------

\vspace{6mm}

\begin{tabular*}{\textwidth}{@{\extracolsep{\fill}}b{90mm}@{\hspace{8mm}}>{\raggedleft\arraybackslash}b{62mm}@{}}
  \begin{minipage}[b]{\linewidth}
    {\fontsize{10.7}{12.8}\selectfont
      \SeminarSeries\par
    }

    \vspace{1mm}

    {\fontsize{10.7}{12.8}\selectfont
      \SeminarVenue\par
    }

    \vspace{1mm}

    {\bfseries\fontsize{10.7}{12.8}\selectfont
      \SeminarTime\par
    }
  \end{minipage}
  &
  \includegraphics[
    trim=35 160 35 155,
    clip,
    width=27mm
  ]{QR.pdf}
\end{tabular*}

\vspace{38mm}

% -----------------------------------------------------------------------------
% TALK TITLE AND SPEAKER
% -----------------------------------------------------------------------------

\begin{center}
  {\bfseries\fontsize{14.5}{16.8}\selectfont
    \SeminarTitle\par
  }

  \vspace{0.8mm}

  {\fontsize{10.7}{12.7}\selectfont
    \textcolor{black}{\SpeakerName} (\SpeakerAffiliation)\par
  }
\end{center}

\vspace{1.3mm}

% -----------------------------------------------------------------------------
% ABSTRACT
% -----------------------------------------------------------------------------

{\fontsize{9.7}{11.55}\selectfont
  \textbf{Abstract.}
  In this talk, I will introduce some results on the linear stability of relative equilibria in the N-body problem obtained using the Maslov-type index. The talk consists of three parts.

In the first part, I will briefly introduce the basic definition of a relative equilibrium and describe the linear stability problem.

In the second part, I will introduce some basic background on the Maslov-type index and explain the idea of using the Maslov-type index to study the linear stability problem.

In the third part, I will present the fundamental solution of the Kepler problem and then some qualitative and quantitative results on the linear stability of relative equilibria in the N-body problem obtained using the idea introduced in the second part.

This talk is based on joint work with Prof. Xijun Hu and Prof. Yuwei Ou.
  \par
}

\vfill
\end{minipage}
\par
\nointerlineskip

% =============================================================================
% LOWER HALF: SPEAKER BIOGRAPHY
% =============================================================================

{\bfseries\fontsize{11}{13}\selectfont
  \SpeakerName\par
}

\vspace{2mm}

{\fontsize{9.2}{11}\selectfont
  \begin{tabularx}{\textwidth}{@{}>{\bfseries}p{25mm}@{\hspace{2mm}}>{\raggedright\arraybackslash}X@{}}
    2026.06--Present
      & PostDoc in Mathematics, Dalian University of Technology, China.\\[0.35mm]
    2026.03--2026.06
      & Visiting Scholar, Nankai University, China.\\[0.35mm]
    2021.09--2025.12
      & PhD in Mathematics, Shandong University, China.\\[0.35mm]
    2023.11--2025.10
      & Visiting PhD student, Paris Observatory, France.\\[0.35mm]
    2018.09--2021.06
      & Master in Mathematics, Shandong University, China.\\[0.35mm]
    2014.09--2018.06
      & Bachelor in Mathematics, Lanzhou University, China.
  \end{tabularx}
  \par
}

\vspace{3mm}

{\fontsize{9.5}{11.3}\selectfont
  \textbf{Research Interests:}
  Relative Eqilibria of the Newtonian N-body Problem, Applications of Maslov Index Theory in the N-body Problem.
  \par
}

\vfill

% =============================================================================
% DLADE BRAND BAR
% =============================================================================

{\color{DLADEGold}\rule{\textwidth}{0.9mm}}\par
\vspace{0.6mm}

\begingroup
\setlength{\fboxsep}{0pt}
\noindent\colorbox{DLADEBlue}{%
  \begin{minipage}[c][10mm][c]{\textwidth}
    \begin{tabular*}{\linewidth}{@{\hspace{6mm}}m{9mm}@{\hspace{2.3mm}}m{70mm}@{\extracolsep{\fill}}>{\raggedleft\arraybackslash}m{60mm}@{\hspace{6mm}}}
      \includegraphics[width=9mm]{dlade-favicon.pdf}
      &
      {\sffamily\bfseries\fontsize{7.2}{8.5}\selectfont\color{white}
        DLADE - DaLian Analysis\,\&\,Differential Equations
      }
      &
      {\sffamily\fontsize{5.8}{7}\selectfont\color{white!78}
        SCHOOL OF MATHEMATICAL SCIENCES - DUT
      }
    \end{tabular*}
  \end{minipage}%
}
\par
\endgroup

\end{document}
